% ~150 words, written LAST. Five sentences: problem, gap, what we did, key result, meaning.
% 1. Damaged historical newspapers give garbled OCR; repair models can invent plausible but wrong facts.
% 2. Gap: no controlled severity scale with a fact-level metric, identical input for encoder and LLM.
% 3. We simulate span damage at 6 severity levels on 150 BLN600 test sentences and compare fine-tuned BERT with dictionary candidates, a few-shot LLM (Qwen3-30B-A3B), a dictionary lookup and no repair (BART descriptively).
% 4. Key result: \todo{fill from test run: where quality breaks down, which method wins, BERTScore vs Fact Recovery gap}.
% 5. Meaning: \todo{how far repair can be trusted as damage grows}.
OCR errors in historical newspapers hit names, places and numbers hardest, and generative repair may replace them
with fluent but wrong words. We damage one span of words around a fact at five nested levels, from one word to 75\%
of the sentence, in 150 \bln{} test sentences, with character noise calibrated on the corpus's own OCR errors. With
identical input, we compare fine-tuned BERT with dictionary candidates, a few-shot LLM (Qwen3-30B-A3B), a dictionary
lookup and, descriptively, a fine-tuned BART denoiser, scoring repairs with BERTScore and an exact Fact Recovery
Rate. We find no single level at which repair breaks down, but no main method restores more than about half of the facts: BERT with
dictionary candidates recovers the anchor fact in 58\% of sentences at one damaged word and 53\% at 75\%, the lookup
alone 51\%. The LLM falls from 49\% to 25\%, partly through format failures, and never beats the lookup on facts.
BERTScore overstates repair: from 10\% damage on, the denoiser scores highest on it while restoring fewer facts.
